\documentclass[12pt,a4paper,oneside]{article}
\usepackage[brazil]{babel}
\usepackage[utf8]{inputenc}
\usepackage{olymp}

\setcounter{problem}{0}

\begin{document}

\begin{problem}{Concurso}{concurso}{2}
 
Em certo concurso, vários candidatos fizeram uma prova de múltipla escolha. O edital do concurso previa que a vaga seria preenchida pelo candidato que obtivesse a maior nota. Mas o estagiário que redigiu o edital se esqueceu de incluir a regra de desempate! Agora que o concurso já foi realizado, os responsáveis precisam saber o tamanho do problema que tem em mãos... descobrir quantas pessoas empataram com a maior nota!

Você deve fazer um programa para ler as notas dos candidatos e informar a maior nota e quantos candidatos obtiveram esta nota.

%\Notes
%
%Observações, se tiver.

\InputFile

A entrada é composta por duas linhas. A primeira contém um valor inteiro $N$, que indica o número de candidatos que fizeram a prova. A segunda linha contém $N$ valores inteiros, indicando a nota de cada
candidato. Restrições: $1 \leq N \leq 1000$ e a nota de cada candidato pode variar de $0$ a $30$.

\OutputFile

Seu programa deve imprimir dois valores separados por espaço em branco: a maior nota e quantos candidatos obtiveram esta nota.

% \Notes
% Preste atenção aos tipos das variáveis, pois $F(90) \approx 2^{61}$.

\Example

\begin{example}
\exmp{
7
20 7 9 13 2 12 23
}{
23 1
}%
\end{example}

\begin{example}
\exmp{
10
7 5 2 4 6 8 9 4 9 3
}{
9 2
}%
\end{example}

\begin{example}
\exmp{
5
5 5 5 5 5
}{
5 5
}%
\end{example}


%Comentários sobre o exemplo, se necessário.

\end{problem}

\end{document}
